\documentclass[11pt,a4paper]{article}

\usepackage[utf8]{inputenc}
\usepackage[T1]{fontenc}
\usepackage[provide=*,english]{babel}
\usepackage{lmodern}
\usepackage{microtype}
\usepackage{amsmath,amssymb,mathtools}
\usepackage{geometry}
\usepackage{xcolor}
\usepackage{enumitem}
\usepackage{fancyhdr}
\usepackage[hidelinks]{hyperref}

\geometry{left=2.3cm,right=2.3cm,top=2.2cm,bottom=2.25cm,headheight=14pt}
\setlength{\parindent}{0pt}
\setlength{\parskip}{0.42em}
\setlist[enumerate]{topsep=0.3em,itemsep=0.25em,leftmargin=1.8em}
\allowdisplaybreaks

\definecolor{blue}{RGB}{34,73,110}
\definecolor{gray}{RGB}{90,96,102}

\pagestyle{fancy}
\fancyhf{}
\fancyhead[L]{\small\color{gray}GC2D: simplified ABBA method}
\fancyhead[R]{\small\color{gray}Projection by averaging}
\fancyfoot[C]{\thepage}
\renewcommand{\headrulewidth}{0.3pt}

\newcommand{\R}{\mathbb{R}}
\newcommand{\Id}{I}
\newcommand{\Y}{\mathcal{Y}}
\newcommand{\M}{\mathcal{M}}
\newcommand{\Om}{\Omega_{\times}}
\newcommand{\PhiAB}{\Phi^{\mathrm{ABBA}}_{h,t_n}}

\begin{document}

\begin{center}
  {\LARGE\bfseries\color{blue}ABBA Method with Projection by Averaging}\par
  \vspace{0.35em}
  {\large Simplified version for the GC2D problem}\par
  \vspace{0.4em}
  {\small No multipliers, no implicit equation, and no Newton iterations}
\end{center}

\vspace{0.45em}

\section{Physical system}

The physical state is denoted by
\[
 y=(X,Y)\in\R^2.
\]
For guiding-center motion, we consider
\[
 \dot X=-\partial_Y\phi(t,X,Y),
 \qquad
 \dot Y=\partial_X\phi(t,X,Y).
\]
Defining
\[
 J_0=
 \begin{pmatrix}
  0&-1\\
  1&0
 \end{pmatrix},
 \qquad
 f(y,t)=J_0\nabla_y\phi(y,t),
\]
the system can be written compactly as
\[
 \dot y=f(y,t).
\]

\section{State duplication}

We introduce two copies of the physical state,
\[
 \Y=(u,v)\in\R^2\times\R^2,
\]
and consider the duplicated dynamics
\[
 \dot u=f(v,t),
 \qquad
 \dot v=f(u,t).
\]
The diagonal
\[
 \M=\{(u,v)\in\R^4:u=v\}
\]
represents the physical space. At the beginning of each step, the known state is duplicated:
\[
 \boxed{u^{(0)}=v^{(0)}=y_n.} \tag{1}
\]

The essential difference from Hairer's symmetric projection is that the two copies are
not separated before applying ABBA. In particular, no multiplier \(\mu\) appears and no
nonlinear system needs to be solved.

\section{Explicit submaps}

For a substep of length \(s\) and an evaluation time \(t\), we define
\[
 A_{s,t}(u,v)=\bigl(u+s f(v,t),v\bigr), \tag{2}
\]
\[
 B_{s,t}(u,v)=\bigl(u,v+s f(u,t)\bigr). \tag{3}
\]
The submap \(A\) modifies only \(u\), using the fixed value of \(v\). The submap \(B\)
modifies only \(v\), using the fixed value of \(u\). Therefore, both submaps are explicit.

We take
\[
 s=\frac h2,
 \qquad
 t_{n+1}=t_n+h,
\]
and use the composition
\[
 \boxed{
 \PhiAB
 =A_{s,t_{n+1}}
  \circ B_{s,t_{n+1}}
  \circ B_{s,t_n}
  \circ A_{s,t_n}.
 } \tag{4}
\]
Compositions are read from right to left. Thus, the effective order is
\[
 A_{s,t_n}\longrightarrow B_{s,t_n}
 \longrightarrow B_{s,t_{n+1}}\longrightarrow A_{s,t_{n+1}}.
\]

\section{One complete step of the method}

Starting from (1), the four ABBA stages are
\begin{align}
 u^{(1)}
 &=u^{(0)}+s f\bigl(v^{(0)},t_n\bigr), \tag{5a}\\
 v^{(1)}
 &=v^{(0)}+s f\bigl(u^{(1)},t_n\bigr), \tag{5b}\\
 v^{(2)}
 &=v^{(1)}+s f\bigl(u^{(1)},t_{n+1}\bigr), \tag{5c}\\
 u^{(2)}
 &=u^{(1)}+s f\bigl(v^{(2)},t_{n+1}\bigr). \tag{5d}
\end{align}
In general, the two copies do not coincide after completing ABBA:
\[
 u^{(2)}\neq v^{(2)}.
\]
The new physical state is defined simply as their arithmetic mean:
\[
 \boxed{
 y_{n+1}=\frac{u^{(2)}+v^{(2)}}{2}.
 } \tag{6}
\]
To begin the next step, the variables are duplicated again:
\[
 \boxed{
 u_{n+1}=v_{n+1}=y_{n+1}.
 } \tag{7}
\]
In this way, every step begins exactly on the diagonal manifold \(\M\).

\section{Interpretation as a projection}

We define the diagonal embedding and mean extraction by
\[
 E=
 \begin{pmatrix}
  \Id_2\\ \Id_2
 \end{pmatrix},
 \qquad
 P=\frac12
 \begin{pmatrix}
  \Id_2&\Id_2
 \end{pmatrix}.
\]
Then
\[
 Ey=(y,y),
 \qquad
 P(u,v)=\frac{u+v}{2},
 \qquad
 PE=\Id_2.
\]
The projection map onto the diagonal is
\[
 \Pi=EP
 =\frac12
 \begin{pmatrix}
  \Id_2&\Id_2\\
  \Id_2&\Id_2
 \end{pmatrix},
\]
and satisfies
\[
 \Pi^2=\Pi,
 \qquad
 \operatorname{Im}(\Pi)=\M.
\]
Therefore, \(\Pi\) is the Euclidean orthogonal projection of \((u,v)\) onto the diagonal:
\[
 \Pi(u,v)
 =\left(\frac{u+v}{2},\frac{u+v}{2}\right).
\]
The complete physical step can be written compactly as
\[
 \boxed{
 y_{n+1}=\Psi_{h,t_n}(y_n),
 \qquad
 \Psi_{h,t_n}=P\,\PhiAB\,E.
 } \tag{8}
\]

\section{Order of the method}

The method retains second order. To verify this, we set
\[
 f_n=f(y_n,t_n),
 \qquad
 W_n=D_yf(y_n,t_n),
 \qquad
 (f_t)_n=\partial_t f(y_n,t_n).
\]
The exact solution satisfies
\[
 y(t_n+h)
 =y_n+h f_n
 +\frac{h^2}{2}\bigl(W_nf_n+(f_t)_n\bigr)
 +\mathcal O(h^3). \tag{9}
\]
Expanding the four stages in (5) gives
\[
 u^{(2)}
 =y_n+h f_n
 +\frac{h^2}{2}\bigl(W_nf_n+(f_t)_n\bigr)
 +\mathcal O(h^3),
\]
\[
 v^{(2)}
 =y_n+h f_n
 +\frac{h^2}{2}\bigl(W_nf_n+(f_t)_n\bigr)
 +\mathcal O(h^3).
\]
Their mean has the same expansion. Consequently, the local error is
\(\mathcal O(h^3)\) and the global error is \(\mathcal O(h^2)\):
\[
 \boxed{\text{the ABBA method with mean projection is second order}.}
\]

\section{Geometric properties and main limitation}

The four submaps \(A\) and \(B\) are symplectic in the duplicated space with respect to
the matrix
\[
 \Om=
 \begin{pmatrix}
  0&J_0\\
  J_0&0
 \end{pmatrix}.
\]
Therefore, the ABBA map \(\PhiAB\) is also symplectic in \(\R^4\). However, the
projection \(\Pi\) has rank two and is not invertible. Consequently, it is not a
symplectic map in the duplicated space, and the symplecticity of \(\PhiAB\) does not imply
the symplecticity of the physical map
\[
 \Psi_{h,t_n}=P\PhiAB E.
\]

In fact, the physical map is not symplectic in general. Consider the autonomous example
\[
 f(y)=J_0y,
\]
corresponding to the quadratic potential
\[
 \phi(X,Y)=\frac12(X^2+Y^2).
\]
In this case, the simplified method produces the linear map
\[
 y_{n+1}=K_hy_n,
 \qquad
 K_h=\left(1-\frac{h^2}{2}\right)\Id_2
     +\left(h-\frac{h^3}{8}\right)J_0.
\]
Its determinant is
\[
 \det K_h
 =\left(1-\frac{h^2}{2}\right)^2
  +\left(h-\frac{h^3}{8}\right)^2
 =1+\frac{h^6}{64}.
\]
For \(h\neq0\), this determinant is not equal to one. Since symplecticity in two
dimensions is equivalent to preservation of oriented area, this example proves that
\[
 \boxed{
 \text{projection by averaging does not guarantee symplecticity.}
 }
\]
The determinant defect in this example is of order \(h^6\), so it may be small for short
steps, but it does not vanish exactly. For the same reason, although the ABBA architecture
is palindromic, the physical map obtained after projection is not exactly reversible in
general either.

\section{Implementation scheme}

One step can be summarized by the following algorithm:
\begin{enumerate}
 \item Receive \(y_n\), \(t_n\), and \(h\).
 \item Set \(s=h/2\) and duplicate the state:
       \(u\leftarrow y_n\), \(v\leftarrow y_n\).
 \item Apply successively:
       \[
       u\leftarrow u+s f(v,t_n),
       \]
       \[
       v\leftarrow v+s f(u,t_n),
       \]
       \[
       v\leftarrow v+s f(u,t_n+h),
       \]
       \[
       u\leftarrow u+s f(v,t_n+h).
       \]
 \item Project by averaging:
       \[
       y_{n+1}\leftarrow\frac{u+v}{2}.
       \]
 \item Use \(u=v=y_{n+1}\) as the input for the next step.
\end{enumerate}

In the nonautonomous case, four evaluations of \(f\) are performed per step. No Jacobian
is computed, no multiplier is introduced, and no Newton iteration is carried out. The
cost of this simplification is the loss of the exact symplecticity and reversibility
guarantees provided by the implicit symmetric projection.

\section{Summary}

The simplified method is defined by three operations:
\[
 \boxed{
 (y_n,y_n)
 \xrightarrow{\ \mathrm{ABBA}\ }
 (u^{(2)},v^{(2)})
 \xrightarrow{\ \mathrm{mean}\ }
 y_{n+1}=\frac{u^{(2)}+v^{(2)}}{2}.
 }
\]
It is a fully explicit, second-order method. Its implementation is much simpler than that
of the method with Hairer's symmetric projection, but the final average is a Euclidean
projection, not a symplectic projection.

\end{document}